\chapter{Basics of counting}
\begin{notation}
Let $A$ be a finite set. We denote $|A|$ the cardinality of $A$, i.e. the number of elements in the set.
\end{notation}
\begin{definition}
Denote by $[n]$ the set of first $n$ natural numbers $[n]:=\{ 1,2,\dots,n \}$
\end{definition}

\section{Bijections}
\begin{theorem}
If there exists a bijection between sets $A$ and $B$ then $|A|=|B|$.
\end{theorem}
\section{Operations with finite Sets}
\begin{exercise}
Suppose that you know $|A|$ and $|B|$. What can you say about $|A\cap B|$, $|A\cup B|$ and $|A\setminus B|$?
\end{exercise}
\begin{definition}[Addition Rule]
If $A\cap B=\varnothing$ then $|A\cup B|=|A|+|B|$
\end{definition}
\begin{definition}[Cartesian Product]
$$
A\times B=\{ (a,b):a\in A,b\in B \}
$$
\end{definition}
\begin{definition}
$$
|A\times B|=|A|\cdot |B|
$$
\end{definition}
\begin{definition}[Disjoint Union]
$$
A\sqcup B=A\times \{ 0 \}\cup B\times \{ 1 \}
$$
\end{definition}
\begin{theorem}
$$
|A\sqcup B|=|A|+|B|
$$
\end{theorem}
\begin{definition}[Exponential object]
$$
A^{B}:= \{ f:f\text{ is a function from }B \text{ to }A \}
$$
\end{definition}
\begin{theorem}
$$
|A^{B}|=|A|^{|B|}
$$
\end{theorem}
\begin{exercise}
What is the number of $n$-letter words in an $m$-letter alphabet?
\end{exercise}
\begin{remark}
All subsets of $A$ are bijective to the number of functions from $A$ to $\{ 0,1 \}$.

Let $B\subset A$ then the function

\begin{align*}
B & \to \mathbb{1_{B}} \\
a & \mapsto \begin{cases}
0 & \text{if } a\not\in B \\
1  & \text{if }a\in B
\end{cases}
\end{align*}

is called the indicative function of $B$.
\end{remark}
\begin{definition}
A permutation $\pi$ of a set $A$ is a bijection $\pi:A\to A$.
\end{definition}
\begin{notation}
The set of permutations of $[n]$ is denoted by $S_{n}$.
\end{notation}
\begin{theorem}
$$
|S_{n}|=n!
$$
\end{theorem}
\begin{exercise}
Construct a bijection between $S_{n}$ and $[1]\times[2]\times\dots \times [n]$.
\end{exercise}
\section{Probabilistic method}
\begin{definition}
A finite probability space is a pair $(\Omega,P)$ where $\Omega$ is a finite set and $P:2^{\Omega}\to[0,1]$ is a function assigning a number from the interval $[0,1]$ to every subset of $\Omega$ such that:
\begin{itemize}[1]
\item $P(\varnothing)=0$
\item $P(\Omega)=1$
\item $P(A\cup B)=P(A)+P(B)$ for any two disjoint sets $A,B\subseteq \Omega$
\end{itemize}
\end{definition}
\subsection{Probability theory to set theory dictionary}

The set $\Omega$ can be thought as the set of all possible outcomes of some random experiment.
\begin{definition}
Elements of $\Omega$ are called elementary events.
\end{definition}
\begin{definition}
Subsets of $\Omega$ are called events.
\end{definition}
\begin{remark}
Elementary events are not events.
\end{remark}
\begin{definition}
Let $w\in \Omega$, $A,B\subseteq \Omega$
\begin{itemize}
\item $w\in A\Leftrightarrow$ event A ocurred
\end{itemize}
\end{definition}
\begin{example}[A random sequence of 0s and 1s]
$\Omega=\{ 0,1 \}^{n}=$ all $n$-term sequences of 0s and 1s with $|\Omega|=2^{n}$.

For $A\subseteq \Omega$ : $P(A):=\frac{|A|}{|\Omega|}=\frac{|A|}{2^{n}}$.

Consider a set $A\subseteq \Omega$ defined as $A:=\{ (w_{1},\dots,w_{n}):w_{1}=1 \}$. $A$ is the event : "the first element of a random sequence is 1".
\begin{problem}
What is the probability of $A$?
\end{problem}
\end{example}

\section{Binomial coefficients}
\begin{definition}
A binomial coefficient $n\choose k$ is the number of ways in which one can choose $k$ objects out of $n$ distinct objects, assuming the order of the elements does not matter.
\end{definition}
\begin{remark}
This is the same as the number of subsets of $k$ elements of an $n$-element set.
\end{remark}
\begin{proposition}
$$
{n\choose k}=\frac{n!}{(n-k)!k!}
$$
\end{proposition}
\begin{notation}
Let $A$ be a finite set and $k$ be a non-negative integer. Then $A\choose k$ is the set of $k$-elements of $A$:
$$
|{A\choose k}|={|A|\choose k}
$$
\end{notation}

\begin{proposition}
The following identities hold:
\begin{itemize}
\item ${n\choose k}+{n \choose k+1}={n+1 \choose k+1}$
\item $n \choose k$ is the $k$-th element in the $n$-th line of Pascal's triangle.
\end{itemize}
\end{proposition}
\begin{proof}
Each subset of the set $[n+1]$ either contains element $n+1$ or not. 

Number of $(k+1)$-element subsets of $[n+1]$ containing $n+1$ is $n\choose k$.

Number of $(k+1)$-elements subsets of $[n+1]$ not containing $n+1$ is $n \choose k+1$.
\end{proof}
\begin{proposition}
The number of subsets of an $n$-element set is $2^{n}$, since we have
$$
2^{n}={n\choose 0}+{n\choose 1}+\dots+{n\choose n}
$$
The number of subsets of an $n$-element set having odd cardinality is $2^{n-1}$. The number of subsets of an $n$-element set having even cardinality is $2^{n-1}$.
\end{proposition}
\begin{proof}
Consider the map $\phi:2^{[n]}\to 2^{[n]}$ defined by
$$
\phi(A)=A\Delta \{ 1 \}=\begin{cases}
A\setminus \{ 1 \} & \text{if }1\in A \\
A\cup \{ 1 \} & \text{otherwise}
\end{cases}
$$
Cardinalities of subsets $A$ and $\phi(A)$ always have different parity.

Since $\phi \circ\phi=\mathrm{Id}$ we deduce that $\phi$ is a bijection between the set of "odd" subsets and the set of "even" subsets.
\end{proof}

\begin{theorem}[Binomial theorem]
$$
(1+x)^{n}={n\choose 0}+{n\choose 1}x+\dots+{n\choose n}x^{n}=\sum_{i=0}^{n} {n\choose i}x^{i}
$$
\end{theorem}
\begin{proposition}
Assume we have $k$ identical objects and $n$ different persons. Then, the number of ways in which one can distribute this $k$ objects among the $n$ persons equals
$$
{n+k-1\choose n-1}={n+k-1\choose k}
$$
\end{proposition}
\begin{remark}
Equivalently, it is a number of solutions to te equation $x_{1}+\dots+x_{n}=k$ in non-negative integers.
\end{remark}
\begin{proof}
Let $A$ be the set of all solutions of equation $x_{1}+\dots+x_{n}=k$ such that $x_{i}\in \mathbf{Z}_{\geq 0}$.

Let $B$ be the set of all subsets of cardinality $n-1$ in $[n+n-1]$.

We construct the bijection $\psi:A\to B$ in the following way
$$
A:= (x_{1},\dots,x_{n})\mapsto B:= \{ x_{1}+1,x_{1}+x_{2}+2,\dots,x_{1}+\dots+x_{n-1}+n-1 \}
$$
We have $x_{1}+x_{2}+\dots+x_{n}=k$. Consider the following string
$$
\underbrace{ \bullet\dots \bullet }_{ x_{1}\text{ bullets} }+\underbrace{ \bullet\dots \bullet }_{ x_{2}\text{ bullets} }+\dots+\underbrace{ \bullet\dots \bullet }_{ x_{n}\text{ bullets} }
$$
This string consists of $k+n-1$ symbols : $k$ "$\bullet$" and $n-1$ "$+$".
$$
\text{Set } B=\{ \text{positions of "+" in the string} \}
$$
We show that $\psi$ is a bijection by computing its inverse map. Let $\mathbf{b}$ be an element of $B$. Suppose that
$$
1\leq b_{1}<b_{2}<\dots<b_{n-1}\leq k+n-1
$$
are the elements of $\mathbf{b}$ written in the increasing order. Then the preimage $\psi ^{-1}(\mathbf{b})$ is an $n$-tuple of integers $(x_{1},\dots,x_{n})$ defined by
\begin{align*}
x_{1} & =b_{1}-1 \\
x_{i} & =b_{i}-b_{i-1},\quad i=2,\dots,n-1 \\
x_{n} & =k+n-1-b_{n-1}
\end{align*}
It is easy to see from these equations that the numbers $x_{i},i=1,\dots,n$ are non-negative integeres and $x_{1}+\dots+x_{n}=k$.

Since there is a bijection between sets $A$ and $B$, their cardinalities are equal adn
$$
|A|=|B|={k+n-1\choose n-1}
$$
\end{proof}

